\documentclass[12pt]{article}
\usepackage[utf8]{inputenc}
\usepackage[T1]{fontenc}
\usepackage[margin=1in]{geometry}
\usepackage{setspace}
\usepackage{xcolor}
\usepackage[colorlinks=true, urlcolor=blue!60!black, linkcolor=blue!60!black]{hyperref}
\setlength{\parindent}{0em}
\setlength{\parskip}{0.6em}
\renewcommand{\baselinestretch}{1.0}
\begin{document}
\begin{center}
{\LARGE \textit{United States v. Butler}, 297 U.S. 1 (1936)\par}
\vspace{0.5em}
{\large Case Brief --- Two-Page Edition\par}
\end{center}

\section{Memory Jogger}
The Spending Clause confers an independent power to tax and spend for the general welfare, but Congress may not use it to coerce or buy regulation of local agriculture (\href{https://supreme.justia.com/cases/federal/us/297/1/}{297 U.S. 1}).

\section{Detailed Case Facts}
The Agricultural Adjustment Act of 1933 paid farmers to cut acreage and production, funded by an earmarked processing tax on the first processing of commodities like cotton. The United States claimed the tax from the receivers of Hoosac Mills, an insolvent Massachusetts cotton mill, including William M.\ Butler.

\section{Procedural History}
The district court ordered payment; the First Circuit reversed; the Supreme Court affirmed.

\section{Judicial Votes}
6--3. Roberts wrote for Hughes, Van Devanter, McReynolds, Sutherland, and Butler; Stone dissented, joined by Brandeis and Cardozo.

\section{Holding}
Processors resisting an earmarked exaction have standing. The general welfare power is not limited to the other enumerated powers (Hamilton, not Madison), but the AAA unconstitutionally regulated agricultural production reserved to the states.

\section{Analysis of Opinions}
Roberts held the exaction part of a regulatory plan, adopted Hamilton's view, and then held that the program coerced farmers, who faced ruin if they declined, and that Congress cannot purchase what it cannot compel. Stone answered that courts judge power, not wisdom; that ``Threat of loss, not hope of gain, is the essence of economic coercion''; and that a spending power without conditions is meaningless. \textit{Steward Machine}, \textit{Helvering}, \textit{Dole}, and \textit{NFIB} later built on both opinions.

\section{Examples: Future Applications}
\textbf{Same-Side:} conditioning all Medicaid funds on a farm-runoff cap is coercive (\textit{NFIB}). \textbf{Opposite-Side:} a new, optional, germane grant is valid. \textbf{Fence-Sitter:} withholding 25\% of existing highway funds.

\section{Critique}
Critics mock Roberts's ``slot-machine'' account of judicial review, note Stone's point that Hamiltonian spending without conditions is incoherent, and call the earmarking distinction formalistic.

\section{Key Quotations}
\begin{quote}
``The power of courts to declare a statute unconstitutional is subject to two guiding principles of decision which ought never to be absent from judicial consciousness. One is that courts are concerned only with the power to enact statutes, not with their wisdom. The other is that, while unconstitutional exercise of power by the executive and legislative branches of the government is subject to judicial restraint, the only check upon our own exercise of power is our own sense of self-restraint. For the removal of unwise laws from the statute books appeal lies not to the courts, but to the ballot and to the processes of democratic government.''
\end{quote}
\noindent --- Justice Stone, dissenting, 297 U.S.\ at 78--79.

\end{document}
